\documentclass[conference]{IEEEtran}
\IEEEoverridecommandlockouts

\usepackage{cite}
\usepackage{amsmath,amssymb,amsfonts}
\usepackage{graphicx}
\usepackage{textcomp} 
\usepackage{xcolor}
\usepackage{listings}
\usepackage{hyperref}
\usepackage{booktabs}
\usepackage{array}

\lstset{
  basicstyle=\ttfamily\footnotesize,
  keywordstyle=\color{blue},
  commentstyle=\color{gray},
  stringstyle=\color{red},
  breaklines=true,
  frame=single,
  numbers=left,
  numberstyle=\tiny\color{gray},
  language=C
}

\def\BibTeX{{\rm B\kern-.05em{\sc i\kern-.025em b}\kern-.08em T\kern-.1667em\lower.7ex\hbox{E}\kern-.125emX}}

\begin{document}

\title{Terminal Snake: An Operating-Systems Case Study in\\
Custom Memory, I/O, and Process Management}

\author{
  \IEEEauthorblockN{Navnit Naman}
  \IEEEauthorblockA{\textit{School of Technology}\\
  \textit{Rishihood University}\\
  Sonipat, India\\
  230085}
  \and
  \IEEEauthorblockN{Kanhaiya Kumar}
  \IEEEauthorblockA{\textit{School of Technology}\\
  \textit{Rishihood University}\\
  Sonipat, India\\
  230062}
}

\maketitle

% ─────────────────────────────────────────────────────────────────────────────
\begin{abstract}
This paper presents \textit{Terminal Snake}, a real-time terminal game written
in ISO C99 that intentionally avoids the C standard library's memory,
string, and math subsystems.  The project serves as a self-contained
demonstration of core operating-systems concepts: a best-fit heap allocator
with bidirectional coalescing, raw-mode terminal I/O via POSIX \texttt{termios},
ANSI/VT100 screen rendering, a cooperative game-loop scheduler, and a
structured error-handling framework.  All five OS-course evaluation domains
(Memory Management, I/O Management, Process Management, Security/Error
Handling, and File-System considerations) are addressed either through
concrete implementation or through deliberate architectural decision and
documented future scope.  The result is a fully playable game with adaptive
difficulty, a consecutive-food score multiplier, a lives system, and a
live memory-utilisation HUD---all built on fewer than 900 lines of hand-crafted
C.  Version history is maintained on GitHub under account \texttt{naman394},
repository \texttt{kanhaiyak23/Snake\_game\_}.
\end{abstract}

\begin{IEEEkeywords}
operating systems, memory management, process management, I/O management,
custom allocator, terminal game, C99, POSIX, best-fit allocation, coalescing
\end{IEEEkeywords}

% ─────────────────────────────────────────────────────────────────────────────
\section{Introduction}

Most hobby games lean on \texttt{malloc}, \texttt{printf}-wrappers, and
framework event loops without considering the subsystems underneath.  This
project inverts that approach: the constraint of \emph{no standard-library
memory, string, or math functions} forces every OS primitive to be made
explicit and auditable.

The five pillars addressed are:

\begin{itemize}
  \item \textbf{Memory Management} --- a 1\,MiB static pool with best-fit
        allocation and bidirectional free-block coalescing.
  \item \textbf{I/O Management} --- raw-mode keyboard capture via
        \texttt{termios} and ANSI/VT100 screen rendering via \texttt{printf}.
  \item \textbf{Process Management} --- a deterministic cooperative game loop
        with \texttt{usleep}-based frame pacing and level-dependent speed
        scaling.
  \item \textbf{Security / Error Handling} --- structured termination codes,
        double-free guards, out-of-bounds checks, and graceful OOM recovery.
  \item \textbf{File Systems} --- persistent high-score storage and
        configuration, deferred to future work with a clear design rationale.
\end{itemize}

% ─────────────────────────────────────────────────────────────────────────────
\section{System Architecture}

\subsection{Module Layout}

The project is divided into two layers, each with strict dependency rules:

\begin{table}[h]
\centering
\caption{Source Module Overview}
\begin{tabular}{lp{5.2cm}}
\toprule
\textbf{Module} & \textbf{Responsibility} \\
\midrule
\texttt{math.c}     & Arithmetic (mul, div, mod, abs, clamp), XorShift32 PRNG \\
\texttt{string.c}   & strlen, strcpy, strcmp, itoa, strtok replacements \\
\texttt{memory.c}   & 1\,MiB static pool, best-fit allocator, coalescing \\
\texttt{keyboard.c} & POSIX termios raw mode, escape-sequence state machine \\
\texttt{screen.c}   & ANSI/VT100 renderer, field-geometry accessors \\
\midrule
\texttt{snake.h/c}  & Game structs, movement, collision, scoring, lives \\
\texttt{main.c}     & Start screen, countdown, game loop, game-over screen \\
\midrule
\texttt{test\_logic.c} & Non-interactive unit tests for library correctness \\
\bottomrule
\end{tabular}
\end{table}

\noindent Library modules are independent of each other; game code may use all
of them.  \texttt{main.c} may not call game-logic internals; it communicates
exclusively through the six-function public API declared in \texttt{snake.h}.

\subsection{Constraint Rationale}

The project prohibits \texttt{<string.h>}, \texttt{<math.h>}, \texttt{malloc},
\texttt{free}, and \texttt{rand}.  The only permitted standard-library symbols
are \texttt{printf}/\texttt{putchar} (terminal I/O) and \texttt{exit}
(process control).  This constraint forces every memory allocation, random
number, and string operation to be implemented and understood from first
principles---precisely the pedagogical intent of an OS course project.

% ─────────────────────────────────────────────────────────────────────────────
\section{Memory Management}

\subsection{Pool Design}

A single \texttt{static char VRAM[1\,048\,576]} array constitutes the entire
heap.  On \texttt{mem\_init()}, it is formatted as one free \texttt{BlockHeader}
spanning the whole array:

\begin{lstlisting}
typedef struct BlockHeader {
    int  size;              /* user payload bytes */
    int  used;              /* 1=allocated, 0=free */
    struct BlockHeader *next;
} BlockHeader;
\end{lstlisting}

All blocks are linked in address order, enabling $O(1)$ adjacency tests needed
for coalescing.

\subsection{Best-Fit Allocation}

On each \texttt{alloc(size)} call, the allocator:
\begin{enumerate}
  \item Rounds \texttt{size} up to the nearest 8-byte boundary (alignment).
  \item Scans the free list for the smallest block that fits (\emph{best-fit}
        policy), breaking early on an exact match.
  \item Splits the chosen block if the remainder exceeds
        \(\text{sizeof(BlockHeader)} + 8\) bytes.
  \item Returns \texttt{NULL} on exhaustion, triggering the OOM error path.
\end{enumerate}

Best-fit was chosen over first-fit because snake segments are all the same
size (\texttt{sizeof(Segment)} = 12 bytes).  Under uniform allocation sizes,
best-fit leaves fewer unusable slivers and is empirically equivalent to
first-fit in fragmentation, while being strictly no worse \cite{knuth1997art}.

\subsection{Bidirectional Coalescing}

On \texttt{dealloc(ptr)}, the allocator:
\begin{enumerate}
  \item Validates the pointer is within \texttt{VRAM} (bounds guard).
  \item Checks \texttt{blk->used} to prevent double-free corruption.
  \item Marks the block free and coalesces \emph{rightward} (absorbs
        consecutive free successors in a \texttt{while} loop).
  \item Finds the predecessor by address-order traversal and coalesces
        \emph{leftward} if the predecessor is free and adjacent.
\end{enumerate}

This two-phase coalescing prevents fragmentation even under the
alloc-one-grow/dealloc-one-shrink pattern of every game tick.

\subsection{Live Memory HUD}

\texttt{mem\_used()} walks the block list and sums used payload bytes.
This value is displayed every frame in the HUD, making heap behaviour
observable during gameplay---an unusual but educationally valuable feature.

% ─────────────────────────────────────────────────────────────────────────────
\section{I/O Management}

\subsection{Keyboard Input}

\texttt{keyboard\_init()} transitions the terminal into raw, non-canonical,
non-echoing mode via POSIX \texttt{termios}:

\begin{lstlisting}
new_tio.c_lflag &= ~(ICANON | ECHO);
new_tio.c_cc[VMIN]  = 0;  /* non-blocking */
new_tio.c_cc[VTIME] = 0;
tcsetattr(STDIN_FILENO, TCSANOW, &new_tio);
fcntl(STDIN_FILENO, F_SETFL, flags | O_NONBLOCK);
\end{lstlisting}

\noindent Setting \texttt{VMIN=0, VTIME=0} with \texttt{O\_NONBLOCK} makes
\texttt{read()} return immediately whether or not a character is waiting.
This is necessary for a real-time game loop that cannot block on user input.

\subsubsection{Escape-Sequence State Machine}

Arrow keys arrive as three-byte ANSI sequences (\texttt{ESC [ A/B/C/D} for
up/down/right/left).  \texttt{keyPressed()} implements a three-state machine
(\texttt{esc\_state} $\in \{0, 1, 2, 3\}$) to reassemble these sequences
across successive \texttt{read()} calls and map them to the WASD key codes
used by game logic.  Stale buffered input is drained between screens to
prevent bleed-through.

\subsubsection{Terminal Restoration}

\texttt{keyboard\_restore()} removes \texttt{O\_NONBLOCK} and calls
\texttt{tcsetattr} with the saved \texttt{old\_tio} snapshot, guaranteeing
the terminal is left in a working state on both normal exit and error paths.

\subsection{Screen Output}

\texttt{screen.c} is a thin wrapper around ANSI/VT100 escape sequences:

\begin{itemize}
  \item \textbf{Alternate screen buffer} (\texttt{\textbackslash{}033[?1049h/l}):
        isolates game rendering from the user's shell history.
  \item \textbf{Cursor control} (\texttt{\textbackslash{}033[row;colH}):
        positions output anywhere on the terminal grid.
  \item \textbf{SGR color codes} (e.g., \texttt{\textbackslash{}033[92m} for
        bright green): style the snake, food, and HUD without external graphics
        libraries.
  \item \textbf{Unicode box-drawing} (\texttt{╔═╗║╚╝}): rendered as raw UTF-8
        bytes via \texttt{screen\_putstr}, replacing the earlier ASCII border.
\end{itemize}

\texttt{screen\_update\_layout()} queries the terminal dimensions via
\texttt{ioctl(TIOCGWINSZ)} on every frame, allowing the game field to scale
dynamically to any terminal size.  All game code accesses field geometry
through accessor functions (\texttt{screen\_field\_x/y/w/h()}) rather than
cached constants, ensuring layout changes propagate without stale reads.

% ─────────────────────────────────────────────────────────────────────────────
\section{Process Management}

\subsection{Game Loop Architecture}

The game runs under a cooperative, single-process loop in \texttt{main.c}:

\begin{enumerate}
  \item \texttt{keyPressed()} --- non-blocking poll for user input.
  \item \texttt{game\_handle\_input()} --- apply direction change or pause/quit.
  \item \texttt{game\_update()} --- advance physics by one tick.
  \item \texttt{game\_render()} --- redraw the entire frame.
  \item \texttt{usleep(game\_speed\_delay())} --- yield CPU for the remainder
        of the frame budget.
\end{enumerate}

This is a textbook \emph{fixed-timestep cooperative loop} \cite{gregory2018game}.
There is no preemptive multitasking; the process voluntarily yields via
\texttt{usleep}, which on macOS/Linux maps to \texttt{nanosleep} at the kernel
level.

\subsection{Frame-Rate Scaling}

\texttt{game\_speed\_delay()} computes the inter-frame sleep duration as a
function of level:

\[
  \text{delay}(\ell) = \text{clamp}(180{,}000 - 15{,}000 \times (\ell - 1),\; 50{,}000,\; 180{,}000) \;\mu\text{s}
\]

Level~1 yields a 180\,ms frame period (5.6\,fps physics), and level~10 yields
50\,ms (20\,fps physics).  All arithmetic uses \texttt{my\_mul} and
\texttt{my\_clamp} from the custom math library; no floating-point operations
are involved.

\subsection{Level Progression}

Level advances every 50 score points:

\[
  \ell = \min\!\left(\left\lfloor \frac{\text{score}}{50} \right\rfloor + 1,\; 10\right)
\]

Because the score multiplier (Section~\ref{sec:multiplier}) accelerates point
gain, high-skill play naturally reaches maximum speed faster.

\subsection{Startup Sequence}

The ordered startup sequence enforces a safe initialisation contract:
\texttt{mem\_init} $\to$ \texttt{keyboard\_init} $\to$ \texttt{screen\_init}
$\to$ \texttt{my\_srand} $\to$ game loop $\to$ \texttt{game\_free} $\to$
\texttt{screen\_cleanup} $\to$ \texttt{keyboard\_restore}.  Resource
acquisition and release are strictly paired, preventing resource leaks across
play-again sessions.

% ─────────────────────────────────────────────────────────────────────────────
\section{Security and Error Handling}
\label{sec:security}

\subsection{Structured Termination Codes}

Game-ending conditions are represented as symbolic constants rather than
magic integers:

\begin{table}[h]
\centering
\caption{Game-End Reason Codes}
\begin{tabular}{ll}
\toprule
\textbf{Constant} & \textbf{Trigger} \\
\midrule
\texttt{GAME\_END\_NONE}      & Game still running \\
\texttt{GAME\_END\_COLLISION} & Wall or self-collision \\
\texttt{GAME\_END\_OOM}       & \texttt{alloc()} returned \texttt{NULL} \\
\texttt{GAME\_END\_QUIT}      & User pressed \texttt{Q} \\
\texttt{GAME\_END\_INTERNAL}  & Linked-list invariant violated \\
\bottomrule
\end{tabular}
\end{table}

The game-over screen reads \texttt{game\_end\_reason} and displays a
specific diagnostic for each case (e.g., ``Memory: alloc() failed --- 1\,MB
heap full or fragmented'' for OOM), giving the user actionable information
rather than a silent crash.

\subsection{Memory Safety}

\begin{itemize}
  \item \textbf{Bounds guard:} \texttt{dealloc()} checks that the pointer lies
        within \texttt{VRAM} before dereferencing; out-of-range pointers are
        silently ignored.
  \item \textbf{Double-free guard:} the \texttt{used} flag is checked before
        marking a block free; a second \texttt{dealloc} on the same pointer is
        a no-op.
  \item \textbf{OOM propagation:} \texttt{alloc()} returning \texttt{NULL} is
        checked at every call site (\texttt{seg\_new}, \texttt{reset\_snake});
        the game transitions to \texttt{GAME\_END\_OOM} rather than
        dereferencing a null pointer.
  \item \textbf{Invariant assertion:} if the tail-search loop in
        \texttt{game\_update()} cannot locate the tail pointer, the tick is
        aborted, the orphaned head segment is freed, and
        \texttt{GAME\_END\_INTERNAL} is set---preventing a use-after-free.
\end{itemize}

\subsection{Input Validation}

\begin{itemize}
  \item \textbf{Boundary check:} before placing the new head, \texttt{my\_in\_bounds()}
        verifies the computed position lies within the field; out-of-bounds
        positions trigger \texttt{lose\_life\_and\_respawn()} rather than
        writing to an invalid segment.
  \item \textbf{Anti-reversal:} \texttt{game\_handle\_input()} blocks the
        180-degree direction change (snake reversing into itself) by checking
        the current velocity before updating \texttt{dx}/\texttt{dy}.
  \item \textbf{Input draining:} \texttt{drain\_input()} flushes the keyboard
        buffer between menu screens, preventing stale escape sequences from
        triggering unintended game actions.
\end{itemize}

\subsection{Terminal Safety}

Both \texttt{keyboard\_restore()} and \texttt{screen\_cleanup()} are always
called on the exit path, ensuring the terminal is never left in raw mode or
on the alternate screen buffer---a common failure mode in terminal games.

% ─────────────────────────────────────────────────────────────────────────────
\section{Gameplay Methodology}

\subsection{Core Tick Loop}

Each call to \texttt{game\_update()} performs, in order:
(1) layout refresh, (2) out-of-bounds food relocation, (3) head position
computation, (4) boundary collision check, (5) self-collision check,
(6) food collision detection, (7) new head allocation, (8) tail removal or
retention, (9) score/level update.  This strict ordering prevents
read-after-write hazards within a single tick.

\subsection{Score Multiplier}
\label{sec:multiplier}

A consecutive-food combo system rewards sustained play without dying.  The
\texttt{combo} counter in \texttt{GameState} increments on each food eaten and
resets to zero on any life loss.  The score awarded per food is:

\[
  \Delta\text{score} = 10 \times \ell \times \min(\text{combo},\; M_{\max})
\]

where $M_{\max} = 8$ is the multiplier cap (\texttt{MAX\_COMBO\_MULT}).  The
HUD highlights the multiplier in bold yellow whenever it exceeds $1\times$,
giving the player real-time streak feedback.

\subsection{Lives System}

The player begins with \texttt{INITIAL\_LIVES} = 3 lives.
\texttt{lose\_life\_and\_respawn()} decrements the counter, resets the combo,
clears all segment memory, rebuilds the snake at field centre, and
repositions all food items.  When lives reach zero, the game ends with
\texttt{GAME\_END\_COLLISION}.  This design separates \emph{session state}
(score, high score, level) from \emph{run state} (snake segments), so the
score persists across respawns within one session.

\subsection{Multi-Food Spawning}

Up to \texttt{MAX\_FOODS} = 3 food items are active simultaneously.
\texttt{place\_food()} uses a rejection-sampling loop: a candidate position
is drawn from \texttt{my\_rand\_range()} and rejected if it overlaps any snake
segment or any other active food item.  Since the snake is never allowed to
fill the field (collision terminates the session first), the loop always
terminates in finite expected time.

% ─────────────────────────────────────────────────────────────────────────────
\section{File System Considerations}

The current implementation holds all game state in memory; no data is
persisted to disk.  This is an intentional deferral, not an oversight.
The project's memory constraint (\texttt{no malloc/free}) was the primary
design driver, and adding file I/O before the allocator was stable would have
introduced a second axis of failure.

A minimal persistence layer would require:
\begin{itemize}
  \item A plain-text or binary record format for the high-score table
        (e.g., a fixed-size struct serialised to \texttt{\textasciitilde/.snake\_score}).
  \item \texttt{open()}/\texttt{read()}/\texttt{write()}/\texttt{close()}
        calls via \texttt{<unistd.h>} (already permitted by the project
        constraint as a POSIX hardware-abstraction exception).
  \item Atomic write-and-rename (\texttt{rename()} over a temporary file) to
        avoid partial writes corrupting the score file on power loss.
\end{itemize}

This is documented as the highest-priority item in the future-scope roadmap
(Section~\ref{sec:future}).

% ─────────────────────────────────────────────────────────────────────────────
\section{Custom Math Library}

\subsection{Multiplication: Shift-and-Add}

\texttt{my\_mul(a, b)} uses the Russian-peasant (binary) method: the larger
magnitude is doubled in each step while the smaller is halved; the larger is
accumulated when the LSB of the smaller is 1.  Time complexity is
$O(\log\min(|a|,|b|))$ bit operations; space is $O(1)$.

\subsection{Division: Recursive Doubling}

\texttt{divide\_positive\_ll(a, b)} halves the problem by doubling the
divisor until it exceeds the dividend, then unwinds the recursion accumulating
the quotient in binary.  Recursion depth is $O(\log(a/b))$.  \texttt{my\_div}
wraps this with sign handling (truncation toward zero).

\subsection{XorShift32 PRNG}

The random number generator uses the XorShift32 algorithm \cite{marsaglia2003xorshift}:

\begin{lstlisting}
rng_state ^= rng_state << 13;
rng_state ^= rng_state >> 17;
rng_state ^= rng_state << 5;
\end{lstlisting}

This produces a full-period sequence of $2^{32}-1$ values with no multiply
in the hot path.  It is seeded from \texttt{time(NULL)} at startup.

% ─────────────────────────────────────────────────────────────────────────────
\section{Results and Metrics}

\begin{table}[h]
\centering
\caption{Implementation Metrics}
\begin{tabular}{lr}
\toprule
\textbf{Metric} & \textbf{Value} \\
\midrule
Total source lines (excl.\ blank/comments) & $\approx 870$ \\
Heap pool size                              & 1\,048\,576 bytes (1\,MiB) \\
Segment size (allocated unit)              & 12 bytes \\
Maximum live segments (level 10, length 200) & $\approx 200$ \\
Peak heap usage at max length              & $\approx 4\,800$ bytes (0.46\%) \\
Frame period range                          & 50\,ms -- 180\,ms \\
Unit tests passing                          & 22 / 22 \\
\bottomrule
\end{tabular}
\end{table}

\noindent The allocator is empirically fragmentation-free for the snake
workload: each tick allocates one segment and (on non-eat frames) frees one
segment of identical size, so the pool converges to a near-steady-state after
a few ticks.

% ─────────────────────────────────────────────────────────────────────────────
\section{Future Scope}
\label{sec:future}

\begin{enumerate}
  \item \textbf{High-Score Persistence (File System):}
        Serialise the top-5 score table to a dot-file using POSIX
        \texttt{open/write/rename} with atomic replace semantics.
  \item \textbf{Slab Allocator:}
        Since all game allocations are the same size (\texttt{sizeof(Segment)}),
        a slab/pool allocator \cite{bonwick1994slab} would reduce per-allocation
        overhead from $O(n)$ (best-fit scan) to $O(1)$ (free-list pop).
  \item \textbf{Configurable Difficulty:}
        Read a plain-text \texttt{.snakerc} configuration file at startup to
        set starting lives, field size, and food count---a direct exercise of
        the deferred file-system module.
  \item \textbf{Replay Recording:}
        Log the RNG seed and input stream to a file; replay by re-seeding
        and feeding the recorded keystrokes, enabling deterministic replay
        without storing frame buffers.
  \item \textbf{Multiplayer over Sockets:}
        Two game instances on the same machine could share state via a
        UNIX-domain socket, exercising IPC and synchronisation primitives.
  \item \textbf{Signal Handling:}
        Install a \texttt{SIGWINCH} handler to respond to terminal resize
        events asynchronously rather than polling \texttt{TIOCGWINSZ} every
        frame, reducing unnecessary system calls.
\end{enumerate}

% ─────────────────────────────────────────────────────────────────────────────
\section{Version History and Collaboration}

All source code, commit history, and issue tracking are maintained at
\texttt{github.com/kanhaiyak23/Snake\_game\_}.  Contributions from both
authors are individually attributable through commit authorship.
Key milestones recorded in the history include:

\begin{itemize}
  \item Initial snake movement and collision (\texttt{e1d954b})
  \item Lives system with respawn (\texttt{fb0e676})
  \item Multi-food spawning and HUD (\texttt{915acd4})
  \item Input-parsing improvements and layout accessors (\texttt{205a8bc})
  \item Score multiplier (combo system) and Unicode border (\texttt{766dafc})
\end{itemize}

% ─────────────────────────────────────────────────────────────────────────────
\section{Conclusion}

Terminal Snake demonstrates that a complete, playable game can be built
atop hand-crafted OS primitives without sacrificing correctness or safety.
The best-fit allocator with bidirectional coalescing correctly manages the
alloc/dealloc pattern of the game loop under all tested conditions.  The
\texttt{termios} raw-mode keyboard driver and ANSI/VT100 renderer together
provide responsive I/O with deterministic terminal restoration.  The structured
error-handling framework ensures that heap exhaustion, list invariant failures,
and user-initiated quits all produce meaningful diagnostics rather than
undefined behaviour.  The project provides a concrete, measurable reference
implementation for each concept covered in the Operating Systems curriculum.

% ─────────────────────────────────────────────────────────────────────────────
\begin{thebibliography}{00}

\bibitem{knuth1997art}
D.~E. Knuth, \textit{The Art of Computer Programming, Vol.~1: Fundamental
Algorithms}, 3rd ed. Reading, MA: Addison-Wesley, 1997, pp.~435--456.

\bibitem{gregory2018game}
J.~Gregory, \textit{Game Engine Architecture}, 3rd ed. Boca Raton, FL: CRC
Press, 2018, pp.~391--420.

\bibitem{bonwick1994slab}
J.~Bonwick, ``The slab allocator: An object-caching kernel memory allocator,''
in \textit{Proc. USENIX Summer Technical Conf.}, Boston, MA, Jun. 1994,
pp.~87--98.

\bibitem{marsaglia2003xorshift}
G.~Marsaglia, ``Xorshift RNGs,'' \textit{J. Statistical Software}, vol.~8,
no.~14, pp.~1--6, 2003.

\bibitem{silberschatz2018os}
A.~Silberschatz, P.~B. Galvin, and G.~Gagne, \textit{Operating System
Concepts}, 10th ed. Hoboken, NJ: Wiley, 2018.

\bibitem{kerrisk2010linux}
M.~Kerrisk, \textit{The Linux Programming Interface}. San Francisco, CA:
No Starch Press, 2010, pp.~1331--1360.

\bibitem{ieee2016posix}
IEEE and The Open Group, \textit{The Open Group Base Specifications Issue 7,
2018 Edition (POSIX.1-2017)}, IEEE Std 1003.1-2017, 2017.

\end{thebibliography}

\end{document}
